\section{Introduction}
\label{sec:introduction}

The dominant format for communicating structured knowledge remains the PDF document~\cite{Lamport1994}. A \LaTeX\ file is a linear stream of characters whose logical structure (sections, theorems, proofs, cross-references) exists only as typesetting conventions. The dependency between Theorem~3.2 and Definition~2.1 has no machine-readable representation; it must be reconstructed by the reader from context.

\begin{figure}[h]
\centering
\begin{tikzpicture}[
  code/.style={draw, rounded corners=2pt, fill=gray!5, font=\ttfamily\scriptsize, text=black, inner sep=4pt, align=left},
  hl/.style={text=red!70!black, font=\ttfamily\scriptsize\bfseries},
  arr/.style={->, >=Stealth, thin, red!60!black}
]
% Definition block
\node[code] (def) at (0,0) {
  \textbackslash begin\{definition\}\textcolor{red!70!black}{\textbackslash label\{def:compact\}}\\
  A subset $K \subseteq \mathbb{R}^n$ is compact if\\
  every open cover has a finite subcover.\\
  \textbackslash end\{definition\}
};
% Prose block
\node[code] (thm) at (7.5,0) {
  Every bounded sequence in $\mathbb{R}^n$\\
  has a convergent subsequence\\
  (by \textcolor{red!70!black}{\textbackslash ref\{def:compact\}}).
};
% Arrow
\draw[arr] ([xshift=-2pt]thm.west) -- ([xshift=2pt]def.east);
\end{tikzpicture}
\caption{A \LaTeX\ cross-reference: \texttt{\textcolor{red!70!black}{\textbackslash label}} and \texttt{\textcolor{red!70!black}{\textbackslash ref}} create a string-level link between two locations. The reference records \emph{that} Theorem~3.2 cites Definition~2.1 but not \emph{why}; the nature of the dependency (used as a hypothesis? as a special case?) must be inferred from surrounding prose.}
\label{fig:latex-ref}
\end{figure}

This is, in a precise sense, the weakest possible structure: a flat string with no first-class morphisms.

Personal knowledge management tools such as Roam Research~\cite{RoamResearch}, Obsidian~\cite{Obsidian}, Logseq~\cite{Logseq}, and Tana~\cite{Tana} introduced bidirectional links, turning notes into graphs. These links, however, lack type information, directional semantics, and composability.

\begin{figure}[h]
\centering
\begin{tikzpicture}[
  note/.style={draw, rounded corners=2pt, fill=gray!5, font=\small, inner sep=4pt, align=left, text width=3.2cm},
  link/.style={->, >=Stealth, thin, gray}
]
\node[note] (ml) at (0,0) {\textbf{Machine Learning}\\[2pt]{\scriptsize A field of AI that uses\\statistical methods to...\\[2pt]\texttt{[[Neural Networks]]}\\[1pt]\texttt{[[Deep Learning]]}\\$\vdots$}};
\node[note] (nn) at (5.5,1.2) {\textbf{Neural Networks}\\[2pt]{\scriptsize Computational models\\inspired by biological...\\[2pt]\texttt{[[Deep Learning]]}\\$\vdots$}};
\node[note] (dl) at (5.5,-1.2) {\textbf{Deep Learning}\\[2pt]{\scriptsize A subset of ML using\\multi-layer networks...\\[2pt]\texttt{[[Machine Learning]]}\\$\vdots$}};
\draw[link] (ml) -- (nn);
\draw[link] (ml) -- (dl);
\draw[link] (nn) -- (dl);
\draw[link, bend left=40] (dl) to (ml);
\end{tikzpicture}
\caption{Obsidian/Logseq: each node is a note file containing prose and \texttt{[[wikilinks]]}. Links are directed (from the note containing the link to the target) but carry no information about the nature of the relationship.}
\label{fig:pkm-links}
\end{figure}

Knowledge graphs in the RDF/OWL tradition~\cite{Bordes2013} address this limitation by introducing typed edges. However, the type system is a fixed enumeration: adding a new relationship type requires modifying the schema itself.

\begin{figure}[h]
\centering
\begin{tikzpicture}[
  dot/.style={circle, fill, inner sep=1.5pt},
  typed/.style={->, >=Stealth, thin}
]
\node[dot] (ml) at (0,0) {};
\node[dot] (nn) at (5,0) {};
\node[dot] (dl) at (2.5,1.3) {};
\node[font=\small, anchor=north] at (ml.south) {Machine Learning};
\node[font=\small, anchor=north] at (nn.south) {Neural Networks};
\node[font=\small, anchor=south] at (dl.north) {Deep Learning};
\draw[typed] (ml) -- node[below, font=\scriptsize\ttfamily] {hasTopic} (nn);
\draw[typed] (ml) -- node[left, font=\scriptsize\ttfamily] {includes} (dl);
\draw[typed] (dl) -- node[right, font=\scriptsize\ttfamily] {isA} (nn);
\end{tikzpicture}
\caption{RDF-style knowledge graph: edges carry a type from a fixed enumeration (\texttt{implies}, \texttt{uses}, \texttt{subsetOf}, \ldots). The type vocabulary is closed; extending it requires schema modification.}
\label{fig:rdf-kg}
\end{figure}

Spivak and Kent~\cite{Spivak2012} proposed \emph{ologs} (ontology logs) as a categorical framework for knowledge representation: objects are types, morphisms are functional relations, and the category-theoretic structure enforces consistency.

\begin{figure}[h]
\centering
\begin{tikzcd}[column sep=large, row sep=large]
\text{a bounded sequence} \arrow[r, "\text{has}"] \arrow[d, "\text{is}"'] & \text{a convergent subsequence} \arrow[d, "\text{is}"] \\
\text{a sequence in } \mathbb{R}^n \arrow[r, "\text{lives in}"'] & \text{a subset of } \mathbb{R}^n
\end{tikzcd}
\caption{An olog~\cite{Spivak2012}: objects are types, morphisms are functional relations, and the diagram commutes. Morphisms are restricted to functions and cannot carry open-ended content.}
\label{fig:olog}
\end{figure}

This was a significant conceptual advance. However, ologs restrict morphisms to functional relations (every object maps to exactly one target) and do not provide a foundation from finite strings, a semantic metric, or an embedding theory. Moreover, the framework has seen limited software adoption in the years since its introduction.

\subsection{Motivation: The Autoformalization Gap}

Our work originates from a concrete software design problem in the formalization of mathematics. In Lean~4~\cite{mathlib4}, the \texttt{leanblueprint} tool~\cite{Massot2020} provides a visual map from informal mathematical statements to their formal counterparts. Dependencies are declared via \LaTeX\ macros: \verb|\uses{lem:foo}| declares that the current statement depends on \texttt{lem:foo}, and \verb|\lean{Bar.baz}| links it to the Lean declaration \texttt{Bar.baz}. These macros generate a dependency graph (Figure~\ref{fig:blueprint-vs-decl}, left), but each edge carries only the label ``uses'' with no further semantic content. LeanArchitect~\cite{Zhu2026} automates blueprint generation, yet the fundamental gap remains: the \emph{meaning} of each dependency is not represented.

\begin{figure}[ht]
\centering
\begin{tikzpicture}[
  bp/.style={draw, rounded corners=2pt, fill=green!10, font=\scriptsize, inner sep=3pt, minimum width=1.4cm, minimum height=0.5cm},
  decl/.style={circle, fill, inner sep=1.5pt},
  arr/.style={->, >=Stealth, thin},
  darr/.style={->, >=Stealth, thin, red!60!black}
]
% Left: Blueprint graph (manual, via \uses{} macro)
\node[font=\small\bfseries] at (1.5, 2.8) {Blueprint};
\node[bp] (la) at (1.5, 2) {\texttt{length\_append}};
\node[bp] (lm) at (0, 0.5) {\texttt{length\_map}};
\node[bp] (mm) at (3, 0.5) {\texttt{map\_map}};
\node[bp] (len) at (1.5, -0.8) {\texttt{List.length}};
\draw[arr] (la) -- node[left, font=\tiny] {uses} (lm);
\draw[arr] (la) -- node[right, font=\tiny] {uses} (len);
\draw[arr] (lm) -- node[left, font=\tiny] {uses} (len);
\draw[arr] (mm) -- node[right, font=\tiny] {uses} (len);

% Right: Declaration-level graph (compiler-extracted, finer-grained)
\node[font=\small\bfseries] at (7.5, 2.8) {Declaration graph};
\node[decl, label={[font=\scriptsize]above:\texttt{length\_append}}] (a) at (7.5, 2) {};
\node[decl, label={[font=\scriptsize]left:\texttt{length\_map}}] (b) at (5.8, 0.8) {};
\node[decl, label={[font=\scriptsize]right:\texttt{map\_map}}] (c) at (9.2, 0.8) {};
\node[decl, label={[font=\scriptsize]below:\texttt{List.length}}] (d) at (7.5, -0.3) {};
\node[decl, label={[font=\scriptsize]left:\texttt{List.rec}}] (rec) at (5.5, -0.5) {};
\node[decl, label={[font=\scriptsize]right:\texttt{Nat.add}}] (add) at (9.5, -0.5) {};
\node[decl, label={[font=\scriptsize]below:\texttt{List.map}}] (map) at (7.5, -1.3) {};
\draw[darr] (a) -- (b);
\draw[darr] (a) -- (d);
\draw[darr] (a) -- (add);
\draw[darr] (b) -- (d);
\draw[darr] (b) -- (map);
\draw[darr] (c) -- (map);
\draw[darr] (c) -- (rec);
\draw[darr] (d) -- (rec);
\draw[darr, dashed, gray] (a) to[bend right=15] (rec);
\draw[darr, dashed, gray] (a) to[bend left=15] (map);
\end{tikzpicture}
\caption{Left: a \texttt{leanblueprint} dependency graph for basic list lemmas. Each edge is labeled ``uses'' via \texttt{\textbackslash uses\{\}} macros. Right: the declaration-level dependency graph extracted from Lean~4 compilation artifacts~\cite{LiPengSeverini2026}. The compiler graph reveals additional nodes (\texttt{List.rec}, \texttt{Nat.add}, \texttt{List.map}) and implicit dependencies (dashed). Both representations lose the semantic content of each edge.}
\label{fig:blueprint-vs-decl}
\end{figure}

Autoformalization~\cite{Wu2022, Jiang2023} aims to translate informal mathematics directly into formal proofs. The gap between ``Draft, Sketch, and Prove'' and full autoformalization is precisely a gap in structured knowledge representation: the system needs to know not just \emph{that} \texttt{length\_append} depends on \texttt{List.length}, but \emph{how}. Does it unfold the definition? Apply it to a subterm? Use it via a rewrite? This information is present in the Lean proof term but lost in both the blueprint and the declaration graph.

In designing Astrolabe to bridge this gap, we realized that the problem is not specific to mathematics. Our prior work on the network structure of Mathlib~\cite{LiPengSeverini2026} made visible the rich structural information latent in formalized mathematics, while our work on content-addressing~\cite{LiParisSeverini2026} demonstrated the advantages of hash-based decentralized identification. Together, these revealed that dependency graphs, even when carefully extracted from compiler artifacts, lose the semantic content of each dependency. The same limitation appears across domains of structured knowledge, including legal codes~\cite{Catala}, software architecture, and scientific literature~\cite{Garfield1955, Price1965}: \emph{morphisms lack open semantics}. Edges in a knowledge graph should carry arbitrary key-value content, just as nodes do.

Every externalized form of human knowledge, whether a \LaTeX\ file, a Lean proof, a Python module, or a legal statute, is a finite string over a finite alphabet. Different parsing procedures (a compiler, an LLM, a human reader) extract different structures from the same underlying string. This observation suggests a universal foundation: start from the free monoid $A^*$ and build upward. A record can freely reference any other record. References are not restricted to ``vertices'' or ``atoms''; an edge can reference another edge, a face can reference a complex. This unconstrained reference structure is the key departure from classical simplicial topology: we impose no constraint on what $\texttt{ref}$ points to. Moreover, a reference list may contain the same vertex multiple times: the \emph{multiplicity profile} $\boldsymbol{\mu}(\sigma)$ records the exact usage pattern and provides the finest combinatorial decomposition of each stage layer.

\subsection{Methodology and Contributions}

We build from a finite alphabet $A$ to the astrolabe file $\mathbb{A} \subset H^+ \times H \times R$ (\S\ref{sec:foundations}), develop its homological algebra via the multicomplex theorem (\S\ref{sec:topology}), add embedding geometry via the stage tower (\S\ref{sec:semantics}), and implement the framework as the open-source tool \texttt{astrolabe.json} (\S\ref{sec:implementation}).

The first contribution is the \emph{Astrolabe} itself: a finite collection $\mathbb{A} \subset H^+ \times H \times R$ of semantic simplices whose reference fields are unconstrained. Each reference list carries a \emph{multiplicity profile} $\boldsymbol{\mu}(\sigma)$ recording how many times each vertex appears. The profile decomposition is the finest partition of each stage layer; degree and stage, the two decompositions that organize classical simplicial topology, are both projections of the profile (\S\ref{sec:foundations}).

The core algebraic result is the \emph{multicomplex theorem} (\S\ref{sec:topology}): grouping restriction maps by which vertex they remove yields $N$ boundary operators $\partial^{(1)}, \ldots, \partial^{(N)}$, one per vertex, that pairwise anticommute. This is strictly stronger than $\partial^2 = 0$: the classical identity is a single linear combination of $\binom{N+1}{2}$ independent equations. Homology, cohomology, long exact sequences, and the snake lemma all follow as corollaries of the anticommutativity, and hold unconditionally.

On the geometric side, LLM embeddings on vertices extend to every stage via a recursive \emph{stage embedding tower} $W_{m+1} = \bigwedge W_m$ (\S\ref{sec:semantics}). The tower is forced, not chosen: stage decomposition requires the next ambient space to contain exterior products of all degrees, which is exactly $\bigwedge W_m$. The topological complexity $k^*$ of a knowledge base gives a provable lower bound $d \geq k^*$ on embedding dimension (the \emph{embedding obstruction theorem}). Grade mixing yields a distance decomposition with an irreducible cross-grade component that has no analogue in classical simplicial geometry.

These tools give formal characterizations of four AI frontier problems: hallucination (grade-stratified: type confusion vs.\ structural flattening), scaling saturation (per-stage: increasing $d$ has exponential payoff at higher stages), the distinction between understanding and memorization (normalized volume: zero means memorizing), and chain-of-thought reasoning (stage-crossing paths exploit dimensional surplus). The characterizations are definitions with testable predictions, not metaphors (\S\ref{sec:semantics}).

On the engineering side, merging two knowledge bases requires no special operation: disjoint union followed by cross-references, with homological invariants tracking the structural effect. Plugins are typed operations (source, endo, sink) composed via the single intermediate representation \texttt{astrolabe.json}. The unified structure reduces pipeline steps from $2p - 1$ to $p$ and annotation cost from $mn$ to $2n$ (\S\ref{sec:implementation}). The algebraic core is mechanically verified in Lean~4 with zero \texttt{sorry}; a reference implementation and the data format are open-source.

\subsection{Related Work}

\begin{itemize}
\item \textbf{Ologs}~\cite{Spivak2012}: No $A^*$ foundation, no restriction maps, no metric, no verification.
\item \textbf{AlgebraicJulia / Catlab.jl}: categorical data structures for scientific modeling. The astrolabe takes a different approach: homological invariants replace categorical universal properties as the primary structural tool.
\item \textbf{CASC}: $\texttt{ref}$ can only reference vertices; no self-enrichment.
\item \textbf{TDA for LLM} (HalluZig, TOHA, G-Detector): Empirical approaches that construct complexes from data. The Astrolabe defines complexes from knowledge structure and explains \emph{why} TDA is effective.
\item \textbf{RDF/OWL}: Edge types are fixed enumerations; no open-ended $\texttt{record}$.
\item \textbf{Obsidian / Logseq}: Edges have no semantics.
\item \textbf{leanblueprint}: Static, manual correspondence between blueprint and formalization.
\item \textbf{$\infty$-categories} (Lurie): Pure theory; no data structure, no implementation.
\end{itemize}
